\section{Acquisition from integrated physical flux}
\label{sec:v11-acquisition}
The preceding construction controls three integer derivatives after
normalizing every observation by $d^2$. We now retain the integrated
flux and use the linear Abel discrepancy of
Section~\ref{sec:v11-abel}. This changes both the approximation loss
and the allocation of preparations. The observation is still exactly
the selected-event bit; no endpoint coordinate or derivative is added.

Keep $\mathcal K_m(B,\beta;D)$, the supplied $g,A,\gamma$, the
labels and \eqref{eq:v10-physical-bounds}. Write
\begin{equation}\label{eq:v11-integrated-observation}
 c_j=2A\sinh(j\gamma),\qquad
 H_{j,b}(d)=c_jp_{j,b}(d)=d^2F_{j,b}(d),\qquad
 H_b=H_{\mathcal V_b}.
\end{equation}
The relative law gives
$\|H_{j,b}-H_b\|_{C^3}\le C\tau^j$. Multiplication by $d^2$
has not altered that estimate. In contrast, the variance of a scaled
empirical integrated flux decreases quadratically at the onset.

\begin{theorem}[Abel-discrepancy acquisition]
\label{thm:v11-acquisition}
Let $L\ge m$, $h=D/L\le1$, $j\ge2$ be even, and
$0<t,\delta\le1$. At every positive node $d_k=kh$ and in each
orientation $b=0,1$, make
\begin{equation}\label{eq:v11-allocation}
 N_{b,k}=\left\lceil
   2c_j(Md_k^2+t/3)t^{-2}\log\frac{4L}{\eta}
                                   \right\rceil
\end{equation}
independent preparations. There is a Borel finite-dictionary estimator
$(\widehat V_0,\widehat V_1)\in\mathcal K^2$ for which, with
probability at least $1-\eta$,
\begin{equation}\label{eq:v11-profile-error}
 \max_b\|\widehat V_b-\mathcal V_b\|_\infty
  \le C\{h^{m-5/2}+\tau^j+t h^{-5/2}+\delta\}.
\end{equation}
The prediction of the odd limiting law has at most $2B$ times this
error. Every preparation, including every failure, is included in the
deterministic bound
\begin{equation}\label{eq:v11-budget-general}
 N_{\rm tot}\le 2L+C A\sinh(j\gamma)h^{-1}t^{-2}
                                       \log\frac{4L}{\eta}.
\end{equation}
No value or derivative at zero is observed. Computational complexity
of the finite dictionary is not included in this preparation bound.
\end{theorem}

\begin{corollary}[A refined preparation bound]
\label{cor:v11-budget}
For sufficiently small $\varepsilon>0$, both complete profiles and
their predicted odd law are recovered to uniform accuracy
$\varepsilon$, with confidence $1-\eta$, using at most
\begin{equation}\label{eq:v11-budget}
 C\varepsilon^{-\left(2+\frac{6}{m-5/2}
                       +\frac{\gamma}{|\log\tau|}\right)}
                         \log\frac{C}{\eta\varepsilon}
\end{equation}
preparations. The flight number is even and of order
$\log(1/\varepsilon)$. This is a sufficient bound for this design,
not a minimax assertion. The factor involving $\tau$ depends on the
supplied relative convergence certificate; it is not asserted to be
an invariant of the billiard.
\end{corollary}

\subsection{A globally smooth positive-node reconstruction}
The Abel discrepancy is applied to globally smooth functions, not to
distributional derivatives across interpolation cells. We give a
specific reconstruction with this property. Let
\[
 \omega(u)=\begin{cases}\exp\{-1/(1-u^2)\},&|u|<1,\\0,&|u|\ge1,
                \end{cases}\qquad
 \theta_\ell(x)=\frac{\omega(x/h-\ell)}
                     {\sum_{r=0}^L\omega(x/h-r)},\quad0\le\ell\le L.
\]
The denominator is positive on $[0,D]$. For each $\ell$, let
$P_\ell z$ interpolate the $m$ positive nodal values starting at
\[
 s_\ell=\max\{1,\min\{\ell-\lfloor(m-1)/2\rfloor,L-m+1\}\}.
\]
Thus the nodes are $s_\ell h,\ldots,(s_\ell+m-1)h$. Define
\begin{equation}\label{eq:v11-smooth-interpolator}
                   J_hz=\sum_{\ell=0}^L\theta_\ell P_\ell z.
\end{equation}
Every polynomial of degree at most $m-1$ is reproduced. At a positive
grid point, only its central weight is nonzero and its polynomial
contains that node, so the nodal values are interpolated. At zero,
the first stencil uses $h,\ldots,mh$; it extrapolates rather than
inserting a known datum. The output is $C^\infty$ on the closed
interval in the usual extension sense.

\begin{lemma}[Smooth reconstruction in the Abel seminorm]
\label{lem:v11-interpolation}
With constants depending only on $m,D$, one has
\begin{align}
 |J_hf-f|_{\mathscr A}
       &\le C h^{m-5/2}\|f\|_{m-1,1},
                    &&f\in C^{m-1,1},\label{eq:v11-approx}\\
 |J_hf|_{\mathscr A}
       &\le C\|f\|_{C^3},
                    &&f\in C^3,\label{eq:v11-smooth-bias}\\
 |J_hz|_{\mathscr A}
       &\le C h^{-5/2}\max_k|z_k|.
                                      \label{eq:v11-scalar-noise}
\end{align}
In particular, a smooth bridge error and an arbitrary nodal error
have different amplification factors.
\end{lemma}
\begin{proof}
At each $x$ at most two weights are positive. A nearest integer to
$x/h$ is at distance at most $1/2$, so the denominator of the weights
has a fixed positive lower bound. Its derivatives of order $r$ are
bounded by $C_rh^{-r}$, as are those of $\theta_\ell$. Each active
stencil lies within distance $(m+1)h$ of $x$. After rescaling, there
are only finitely many Lagrange basis configurations; their derivatives
of every order at most $m-1$ on the relevant compact intervals are
uniformly bounded. The product rule therefore gives
\[
 \|J_hz\|_{C^r}\le C h^{-r}\max_k|z_k|,\qquad 0\le r\le3.
\]
For $f\in C^{m-1,1}$, subtract its degree-$(m-1)$ Taylor polynomial
at the point at which a derivative is evaluated. The interpolator
reproduces that polynomial; its nodal remainders are bounded by
$C h^m\|f\|_{m-1,1}$. The differentiated remainder estimate and
the product rule give
\begin{equation}\label{eq:v11-interpolation-r}
 \|J_hf-f\|_{C^r}\le C h^{m-r}\|f\|_{m-1,1},\qquad r=2,3.
\end{equation}
This argument is valid at both endpoints because all active stencils
remain a bounded rescaled distance away. Subtracting instead the
quadratic Taylor polynomial for $f\in C^3$ gives
$\|J_hf\|_{C^3}\le C\|f\|_{C^3}$.

Apply \eqref{eq:v11-mixed-bound}. In \eqref{eq:v11-interpolation-r}
the product of the second- and third-derivative errors is of order
$h^{2m-5}$; its square root is $h^{m-5/2}$. The endpoint
second-derivative error is $O(h^{m-2})$, a smaller bound when
$h\le1$. For nodal noise the two derivative scales are $h^{-2}$
and $h^{-3}$, whose geometric mean is $h^{-5/2}$. Finally the
$C^3$ stability just proved bounds the smooth-bias term without any
negative power of $h$.
\end{proof}

\subsection{The estimator and its full preparation charge}
The maps $V\mapsto H_V$ take $\mathcal K$ into a bounded
$C^{m-1,1}$ set. This follows from
Lemma~\ref{lem:v10-dictionary} and multiplication by $d^2$.
For fixed $h$, the map
$V\mapsto\mathscr A J_hH_V$ has compact image in $C^0$:
it is a continuous map of the compact class $\mathcal K$ through a
finite-dimensional vector of nodal values. Choose, independently of
the unknown profile, a finite ordered set
$\mathcal D^{\mathscr A}_{h,\delta}\subset\mathcal K$ such that
\begin{equation}\label{eq:v11-dictionary}
 \inf_{W\in\mathcal D^{\mathscr A}_{h,\delta}}
               |J_hH_V-J_hH_W|_{\mathscr A}\le\delta
                       \quad(V\in\mathcal K).
\end{equation}
For a data vector $z$, let $\widehat V(z)$ be the first minimizer of
$|J_hH_W-J_hz|_{\mathscr A}$ over this ordered set. Each criterion
is continuous in $z$; the fixed tie rule makes the estimator Borel.
The affine nullspace of $\mathscr A$ causes no problem for this
finite minimization or for the norm on actual profile differences.

\begin{proof}[Proof of Theorem~\ref{thm:v11-acquisition}]
Let $Z_{b,k}=c_j\overline Y_{b,k}$, with $\overline Y$ the empirical
success proportion. Its mean is $H_{j,b}(d_k)$, and
\[
 \operatorname{Var}(Z_{b,k})
    =\frac{c_jH_{j,b}(d_k)(1-p_{j,b}(d_k))}{N_{b,k}}
    \le\frac{c_jMd_k^2}{N_{b,k}}.
\]
The centered summands are bounded by $c_j$. The elementary Bernoulli
exponential-moment bound used in Section~\ref{sec:v10-acquisition}
therefore gives
\[
 \Pr\{|Z_{b,k}-H_{j,b}(d_k)|>t\}
 \le2\exp\left\{-\frac{N_{b,k}t^2}
                         {2c_j(Md_k^2+t/3)}\right\}.
\]
The allocation and a union bound over $2L$ means give simultaneous
nodal error at most $t$ with probability $1-\eta$.

On this event, Lemma~\ref{lem:v11-interpolation} and the relative
$C^3$ bound give
\[
 |J_hZ_b-J_hH_b|_{\mathscr A}
                    \le C(t h^{-5/2}+\tau^j).
\]
The minimizing property and \eqref{eq:v11-dictionary} bound
$|J_hH_{\widehat V_b}-J_hH_b|_{\mathscr A}$ by twice this
quantity plus $\delta$. The interpolation error for the two
candidate fluxes then gives
\[
 |H_{\widehat V_b}-H_b|_{\mathscr A}
       \le C(h^{m-5/2}+t h^{-5/2}+\tau^j+\delta).
\]
Here the triangle inequality is that of the linear Abel seminorm.
The inverse estimate \eqref{eq:v11-two-sided} proves
\eqref{eq:v11-profile-error}. The positive mass-one bilinear kernel
of $T$ gives the stated odd-law error.

Finally $\sum_{k=1}^L d_k^2\le D^2L$ and $L=D/h$.
Summing \eqref{eq:v11-allocation}, including each ceiling, yields
\eqref{eq:v11-budget-general}. This sums the allocations themselves,
not the number of successes or a waiting time conditional on success.
\end{proof}

\begin{proof}[Proof of Corollary~\ref{cor:v11-budget}]
Put $\nu=m-5/2>0$. Choose $L$ so that $h$ is comparable to a
sufficiently small multiple of $\varepsilon^{1/\nu}$, take
$t=c\varepsilon h^{5/2}$ and $\delta=c\varepsilon$, and choose
the least even $j\ge2$ with $\tau^j\le c\varepsilon$.
All constants can be chosen to control also the odd-law error.
Rounding $j$ changes it by less than two once the lower endpoint is
inactive, so $\sinh(j\gamma)\le
C\varepsilon^{-\gamma/|\log\tau|}$. Moreover
\[
 h^{-1}t^{-2}=c^{-2}\varepsilon^{-2}h^{-6}.
\]
The logarithmic factor is bounded by
$C\log\{C/(\eta\varepsilon)\}$, and the ceiling contribution
$2L$ is smaller than the displayed bound. This proves the claim.
\end{proof}

There are two separate improvements over
Corollary~\ref{cor:v10-budget}: the approximation/noise loss is
$5/2$ rather than three, and the sum of integrated-flux variances has
order $h^{-1}$ rather than the $h^{-2}$ sum produced by dividing
first by $d_k^2$. Neither improvement follows by merely substituting
a mixed-norm inequality in the earlier discrepancy criterion. In
particular, when $m=4$ the accuracy exponent becomes
$6+\gamma/|\log\tau|$ rather than
$10+\gamma/|\log\tau|$. Both remain sufficient exponents with a
supplied convergence certificate.
